%Root main.tex
%---------------------------------------------------------------------------
%第1章　研究背景・目的
%---------------------------------------------------------------------------

\chapter{質疑応答}
\label{chap:first}
\begin{enumerate}
\item[Q1．]補正値を加えていないシミュレーション結果と補正値を加えたシミュレーション結果の波形を比較したとき，どの波形を比較したのか(石川教授)
\item[A1．]補正値を加えていないシミュレーション結果と補正値を加えたシミュレーション結果の出力電流$I_o$を比較した
\item[Q2．]補正値を加えていないシミュレーション結果と補正値を加えたシミュレーション結果のどの結果を評価して平滑化されたといえるのか(石川教授)
\item[A2．]補正値を加えていないシミュレーション結果と補正値を加えたシミュレーション結果の出力電流の波形と全高調波歪み率を評価し，波形がなめらかになり，THDが減少したことから平滑化されたといえる
\item[Q3．]なぜ指令値電流に補正値を加えることで出力電流が改善すると考えたのか(石川教授)
\item[A3．]フィルタリアクトルに流れる電流に発生している高調波成分が除去しきれないため，フィルタリアクトルに流れる電流を恣意的にゆがませて改善を行えると考えたから
\item[Q4．]なぜ先行研究で高調波成分が考慮されていなかったのか（学生）
\item[A4．]先行研究で連系リアクトルを含まない制御モデルに対してデッドビート制御を行ったときに発生したものであるため
\end{enumerate}